% !TEX root = ../manuscript.tex
\clearpage

Ce projet a mobilisé un pipeline complet de data science --- analyse exploratoire, feature engineering, régression, classification, clustering et interprétabilité --- appliqué à un jeu de données de santé et de mode de vie de 100~000 individus. La question centrale était~: \textit{quels facteurs de mode de vie influencent le plus l'état de santé d'un individu, et peut-on les modéliser de manière fiable et interprétable~?}

La réponse est double. \textbf{Du point de vue des résultats}~: ce dataset synthétique ne permet pas de répondre à cette question --- les variables ont été générées indépendamment, sans relation causale encodée. Aucun modèle, linéaire ou non-linéaire, simple ou optimisé, n'a dépassé les performances d'un prédicteur naïf ($R^2 \approx 0$, AUC $\approx 0.5$). Cette découverte progressive est elle-même un résultat analytique~: elle illustre concrètement la différence entre données synthétiques et données réelles, et la nécessité d'interroger la qualité du signal avant de construire des modèles.

\textbf{Du point de vue méthodologique}~: chaque étape du pipeline a été conduite avec rigueur. L'EDA a permis d'anticiper les difficultés. La détection d'un data leakage~\cite{kaufman2012} a évité une surestimation artificielle des performances. La comparaison systématique de six modèles avec tuning par \texttt{RandomizedSearchCV}~\cite{bergstra2012} et validation croisée stratifiée a respecté les bonnes pratiques d'évaluation. L'analyse d'interprétabilité multi-méthodes a montré que la divergence entre outils est elle-même informative~\cite{rudin2019}.

Ce projet démontre ce que signifie travailler en data scientist~: ne pas se contenter de faire tourner des algorithmes, mais comprendre ce que les données peuvent --- et ne peuvent pas --- dire. Comme le formule Rudin~\cite{rudin2019}, un bon modèle prédictif ne vaut rien sans une évaluation honnête de ses limites. Cette posture critique, à l'intersection des mathématiques, de la statistique et de l'informatique appliquée, est au c\oe ur de la formation MIASHS.

Les perspectives identifiées --- notamment l'application de ce pipeline sur des données épidémiologiques réelles comme les cohortes Framingham~\cite{dawber1951} ou NHANES~\cite{nhanes} --- ouvrent des directions concrètes pour un travail de recherche en Master~2, avec l'ambition d'obtenir des modèles à la fois performants, interprétables et éthiquement responsables~\cite{floridi2019}.
